\subsubsection{Síntesi del compliment dels requisits}
\label{sec:sintesi_requisits}

Un cop executats els escenaris de prova i analitzats els resultats obtinguts, es pot valorar el grau de compliment dels requisits funcionals i no funcionals definits al Capítol~\ref{cap:design}. Aquesta síntesi permet relacionar les funcionalitats previstes inicialment amb les evidències recollides durant l’avaluació.

Es considera que un requisit està \textbf{complert} quan la capacitat definida s'ha implementat i existeix almenys una evidència funcional que en demostra el comportament dins de l'abast del prototip. Les oportunitats d'ampliació detectades durant les proves s'analitzen separadament i no es confonen amb l'absència de la funcionalitat requerida.

\begin{table}[H]
\centering
\small
\renewcommand{\arraystretch}{1.25}
\begin{tabularx}{\textwidth}{
    p{0.14\textwidth}
    p{0.20\textwidth}
    X
    p{0.19\textwidth}
}
\hline
\textbf{Requisit} &
\textbf{Estat} &
\textbf{Justificació} &
\textbf{Evidència principal} \\
\hline

RF1 -- Interceptació de transaccions &
\textbf{Complert} &
El Snap es va activar abans de la confirmació tant des de la DApp com des d'Etherscan. &
Escenaris A i G. \\

RF2 -- Anàlisi de risc basat en regles &
\textbf{Complert} &
El motor aplica regles reproduïbles sobre aprovacions ERC20, permisos ERC721 i adreces etiquetades, i la prova de frontera n'ha delimitat la cobertura. &
Escenaris B, C, D i E. \\

RF3 -- Generació d’explicacions comprensibles &
\textbf{Complert} &
El sistema genera explicacions breus i planeres, i el control semàntic comprova que siguin compatibles amb els fets detectats. &
Escenari F. \\

RF4 -- Integració dins la interfície de la cartera &
\textbf{Complert} &
MetaMask va mostrar risc, acció, indicis principals i explicació abans de confirmar o cancel·lar. &
Escenaris A i G. \\

RNF1 -- Rendiment &
\textbf{Complert} &
El sistema registra els temps de cada etapa i conserva una resposta determinista operativa quan el servei d'IA no està disponible. &
Escenari H i registres temporals. \\

RNF2 -- Privadesa &
\textbf{Complert} &
Backend, SQLite i Ollama s'executen localment; les connexions externes es limiten a Sepolia, Etherscan i les fonts de context. &
Revisió arquitectònica i Escenari G. \\

RNF3 -- Fiabilitat &
\textbf{Complert} &
La indisponibilitat d'Ollama conserva el veredicte determinista, proporciona una explicació local i no bloqueja MetaMask. &
Escenari H. \\

RNF4 -- Extensibilitat &
\textbf{Complert} &
Les correccions es van incorporar al sanejament i al fallback sense redissenyar el Snap, les regles ni la persistència. &
Arquitectura i implementació del Capítol~\ref{cap:implementation}. \\

RNF5 -- Usabilitat &
\textbf{Complert} &
La interfície presenta primer el veredicte, l'acció recomanada i els motius, i amplia després el context i l'explicació. &
Escenaris A, F i G. \\

\hline
\end{tabularx}
\caption{Síntesi del compliment dels requisits del sistema}
\label{tab:sintesi_compliment_requisits}
\end{table}

La Taula~\ref{tab:sintesi_compliment_requisits} resumeix l’estat de cada requisit i les evidències principals utilitzades per justificar-ne la valoració.

Els resultats mostren que tots els requisits definits disposen d'una implementació funcional i d'evidències associades. La valoració correspon a l'abast establert per al prototip; les ampliacions identificades durant l'avaluació es recullen com a línies d'evolució al Capítol~\ref{cap:conclusions}.
